\section{Certificates in degrees below 150}\label{sec:certificates}

We now combine the upper bound \(F\le U\), exact target multiplicities,
and highest-weight multiplication. The base interval \(6\le b\le25\)
is checked directly. Beyond it, exact target values supply lower-bound
seeds: they are computed from the target recurrence without assuming
the desired comparison. The resulting certificates leave a set of
partitions whose exact differences will be checked in
Section~\ref{sec:band}. We retain the full interval through \(b=149\)
for the independent chart proof; only \(b\le29\) is needed with
quadratic stabilization.

\subsection{Initial tables and the base interval}
Compute \(Q\) using \eqref{eq:targetrec} on the following four sets.
All partitions have length at most six.
\begin{equation}\label{eq:tables}
\begin{array}{c|l|c}
\text{degrees}&\text{additional restriction}&\text{variables used}\\ \hline
0\le b\le25&\text{none}&6\\
0\le b\le149&\operatorname{length}(\lambda)\le3&3\\
0\le b\le55&\operatorname{length}(\lambda)\le4&4\\
0\le b\le45&t_3(d)\le18&6
\end{array}
\end{equation}
Previous-degree lookups remain in the applicable set by
Section~\ref{sec:upper}. The first table supplies the direct base
comparisons; all four supply lower-bound seeds for larger degrees.
Overlaps are immaterial.

For \(6\le b\le25\), every partition is first tested against \(Q(d)\ge
U(d)\). When this fails, compute \(F(d)=H_6(\lambda(d))\) exactly by
\eqref{eq:sourcerec}. Only these requested values at the last stage
\(j=6\) are needed. Table~\ref{tab:base} records all the fallback
comparisons and their equalities.

\begin{table}[ht]
\centering
\begin{tabular}{r|r|r@{\qquad}r|r|r}
\(b\)&exact comparisons&equalities&\(b\)&exact comparisons&equalities\\ \hline
6&2428&2428&16&7370&354\\
7&4476&315&17&7599&369\\
8&6820&255&18&7801&384\\
9&6899&250&19&8057&399\\
10&6886&264&20&8233&414\\
11&6814&279&21&8427&429\\
12&6850&294&22&8650&444\\
13&6872&309&23&8850&459\\
14&6973&324&24&9031&474\\
15&7131&339&25&9268&489
\end{tabular}
\caption{The base interval: \(145435\) exact fallback comparisons,
\(9272\) equalities, and no negative difference. All other partitions
pass \(Q\ge U\).}\label{tab:base}
\end{table}

\subsection{Residues and transport of lower bounds}
For \(26\le b\le149\) we need only consider \(d_6<6\), by
Lemma~\ref{lem:rectangles}. There is a unique expression
\[
 d=r+6m,\qquad
 r\in\mathcal R=\{0,\ldots,5\}^6\cap\{6\mid w(r)\},\qquad
 m\in\mathcal G=\{m\in\NN^6:m_6=0,\ w(m)\le149\}.
\]
Write \(b_0(r)=w(r)/6\), so that \(b=b_0(r)+w(m)\).
There are \(7776\) residues: choose \(r_2,\ldots,r_6\), and \(r_1\) is
uniquely determined modulo six. Moreover \(b_0(r)\le17\).
The grid \(\mathcal G\) has \(6611697\) points. We call its indices
\(m\) modes: fixing \(m\) groups the \(7776\) possible residues for
simultaneous testing. Pairs whose degree lies outside
\(26\le b\le149\) will not require a test.

For an array on a downward-closed set, \emph{propagation} replaces its
value at \(m\) by the maximum of all values at indices \(u\le m\).
Extending an array means filling new positions by zero and propagating.
An increasing-\(w\) sweep taking maxima with the five immediate
predecessors performs this operation. A reverse sweep through the same
predecessors forms the downward closure of a set.
Lemma~\ref{lem:rectangles} shows that propagation preserves every lower
bound for \(Q(r+6m)\), with \(r\) fixed.

The same multiplication lemma also transports a multiplicity from
residue zero to residue \(r\). Suppose \(J(u)\le Q(6u)\) for
\(u\in\mathcal G\), \(M(m)\le Q(r+6m)\), and \(E\subseteq\mathcal G\)
is a set of indices satisfying
\(Q(r+6c)>0\) for \(c\in E\). Multiplication by one nonzero
highest-weight polynomial at \(r+6c\) gives
\begin{equation}\label{eq:transport}
 Q(r+6m)\ge
 \max\left(\{M(m)\}\cup
      \{J(m-c):c\in E,\ c\le m\}\right).
\end{equation}
The right side certifies the comparison whenever it is at least
\(U(r+6m)\). No product of two multiplicity dimensions is used: a
single nonzero factor gives each injection.

We first construct \(J\) and choices of \(M,E\) valid for every residue,
so that one test can discard a whole mode. For the modes left over we
use more precise choices depending on \(r\). The following construction
specifies all arrays and tests; the common cap \(Z=2^{100}\) is applied
only to lower bounds.

\subsection{Certificates shared by all residues}

\paragraph{A shared power table.}
Initialize \(J(m)=1\) on \(\mathcal G\), using the rectangle occurrences
and the nonzero constant at \(m=0\). For every gap vector \(x\) with
\(x_6=0\), \(Q(x)>1\), and known value in one of the first three sets
of \eqref{eq:tables}, put \(d_0=w(x)/6\) and \(C=\min(Q(x),Z)\).
For each integer \(k\ge1\) with \(kd_0\le149\) and all entries of \(kx\)
divisible by six, update
\[
             J(kx/6)\ \gets\
                 \max\{J(kx/6),\ \min(Z,f(C,z,k))\},
\]
where \(z\) is as in Proposition~\ref{prop:power}. Then propagate.
It follows that
\begin{equation}\label{eq:Jbound}
                         J(m)\le Q(6m)\quad(m\in\mathcal G).
\end{equation}

\paragraph{Common multiplicities and occurrences.}
For each \(r\), initialize the exact array \(Q(r+6u)\) on
\(w(u)\le25-b_0(r)\), then extend it to \(w(m)\le42\).
Take the pointwise minimum of these propagated arrays over all \(r\),
with a cap at \(Z\), to obtain \(M_0\) on that grid.
Let \(E_0\) be the minimal points where \(M_0>0\), before extending
\(M_0\) to all of \(\mathcal G\). Then
\begin{equation}\label{eq:Mbound}
 M_0(m)\le Q(r+6m)\quad\text{for every }r,\qquad
 Q(r+6c)>0\quad\text{for every }r\text{ and }c\in E_0.
\end{equation}
The exact entries supplying the minimum may come from different
predecessors for different \(r\). Their common consequence is precisely
the two universal inequalities in \eqref{eq:Mbound}.

\paragraph{Selecting modes.}
Keep a mode \(m\in\mathcal G\) unless either \(w(m)+17\le25\), or
\begin{equation}\label{eq:sharedtest}
 \max\left(\{M_0(m)\}\cup
       \{J(m-c):c\in E_0,\ c\le m\}\right)
                   \ge U(6m+(5,5,5,5,5,5)).
\end{equation}
The first condition puts every residue in the already settled base
range whenever its degree is at least six. For the second, use
\eqref{eq:transport} with \(M=M_0\) and \(E=E_0\).
Proposition~\ref{prop:upper} bounds \(F(r+6m)\) by
\(U(6m+(5,5,5,5,5,5))\), because \(r\le(5,5,5,5,5,5)\).
The high-residue gap vector need not be admissible; \(U\) is defined and
monotone there nonetheless.

\subsection{Tests for each remaining residue}
Form the downward closure of the retained modes. On this closure,
initialize \(Q(r+6m)\) wherever it is available in
\eqref{eq:tables}, use zero elsewhere, and propagate to obtain \(M_r\).
Let \(E_r\) be the minimal positive locations of the exact array
\(Q(r+6u)\) on \(w(u)\le25-b_0(r)\).
For each originally retained mode with \(26\le b_0(r)+w(m)\le149\),
put \(d=r+6m\) and clear the pair if either
\[
 t_1(d)\le b
 \quad\text{or}\quad
 \max\left(\{M_r(m)\}\cup
          \{J(m-c):c\in E_r,\ c\le m\}\right)\ge U(d).
\]
Flag it otherwise. The first test is Corollary~\ref{cor:first-row}.
The second is \eqref{eq:transport} with \(M=M_r\) and \(E=E_r\).

Minimal positive locations can be found by checking the immediate
predecessors. Indeed positive exact multiplicities form an upper set
within the small grid, by the rectangular shifts; the propagated common
array has the same property. The set used for \(M_r\) retains every
predecessor needed for propagation. Each difference \(m-c\) in a mixed
test is nonnegative and lies in \(\mathcal G\).

In Appendix~\ref{app:certificate}, \texttt{powerAll} stores \(J\),
\texttt{uniBase} and \texttt{uOcc} store \(M_0\) and \(E_0\).
The routine \texttt{setup} retains the modes and their downward closure,
\texttt{init\_occ} constructs \(E_r\), and \texttt{TestSpec} propagates
\(M_r\) and tests the remaining pairs. The routine \texttt{run} performs
the base-interval comparisons.

\begin{verification}\label{ver:certificates}
The complete program in Appendix~\ref{app:certificate} verifies the
base comparisons in Table~\ref{tab:base}. Its mode list has \(87354\)
points, or \(114314\) after taking the downward closure. It flags
\(4969698\) pairs in the middle interval, all with \(t_2(d)\le27\).
Their counts by degree range appear in Table~\ref{tab:band}.
\end{verification}

Every pair in \(26\le b\le149\) cleared by these tests satisfies
\(Q\ge F\), by the proved upper and lower bounds or the chart.
Verification~\ref{ver:certificates}
supplies the separate finite assertion that every pair left over has
\(t_2\le27\). Thus it remains only to check the entire band
\(t_2\le27\) in \(26\le b\le149\).
